% ===== 原文 pp.31–32 =====
% tr-05 片段 a：ψ 类与相交数的定义、Witten 配分函数、string 与 dilaton 方程（p.31–32）
% 起点：p.30/p.31 跨页句 "…More generally using exponents…define" 由 tr-04 写完，
%       本片段从紧随其后的 ⟨τ_0^{d_0}…⟩_g 公式开始。
% 本块无编号公式；无编号陈述；无图。

% [原文 p.31]
\[
\langle\tau_0^{d_0}\tau_1^{d_1}\cdots\tau_m^{d_m}\rangle_g
=\int_{\Mbar_{g,d}}\prod_{j=0}^{m}\prod_{k=1}^{d_j}\psi_{(j,k)}^{j}.
\]

（$\tau_j^{d_j}$ 表示：对 $d_j$ 个穿刺点，相应的 $\psi$ 取 $j$ 次幂；下标 $(j,k)$ 是不同的
穿刺点标号。）对第二种配对，穿刺点的形式计数为 $d=\sum_{j=0}^{m}d_j$，乘积的形式次数是
$\psi$ 因子的计数 $\sum_{j=0}^{m}j d_j$。对非平凡的配对，亏格、穿刺点数与次数满足
$3g-3+d=\sum_{j=0}^{m}j d_j$，否则该配对定义为零。

% [Tutor Note: p.31 的“非平凡配对”判据 $3g-3+d=\sum_{j=0}^{m}j d_j$ 是维数计数：$\dim_\C\Mgn=3g-3+n$（此处 $n=d$），而 $\prod_{j=0}^{m}\prod_{k=1}^{d_j}\psi_{(j,k)}^{j}$ 的总次数为 $\sum_{j=0}^{m}j d_j$，两者不等时按本段约定该相交数取零。代入 $(g,n)=(0,3)$、$d_0=3$ 得 $0=0$，正对应 p.32 的归一化 $\langle\tau_0^3\rangle_0=1$；代入 $(g,n)=(1,1)$、$d_1=1$ 得 $3-3+1=1=\sum_j j d_j$，即 $\langle\tau_1\rangle_1$ 是非零的基本例子。]

已知 psi 类是正的——乘积在子簇上的积分是正的；非平凡配对是正的，这与 Mirzakhani 的
体积多项式系数的正性相符，见定理 4。

\begin{sloppypar}
Witten 考虑了一个配分函数（状态的概率），用于二维引力。对无限向量
$\mathbf t=(t_0,t_1,\dots,t_n,\dots)$，以及 $\gamma$ 的形式求和
$\gamma=\sum_{j=0}^{\infty}t_j\tau_j$，Witten 引入了 $\tau$ 乘积的亏格 $g$ 生成函数
\end{sloppypar}
\[
F_g(\mathbf t)=\sum_{n=0}^{\infty}\frac{\langle\gamma^n\rangle_g}{n!},
\]
其中分子由单项式展开定义，由此得到形式幂级数
\[
F_g(\mathbf t)=\sum_{\{d_j\}}\Big\langle\prod_{j=0}^{\infty}\tau_j^{d_j}\Big\rangle_g
\prod_{j=0}^{\infty}\frac{t_j^{d_j}}{d_j!},
\]
其中求和取遍所有只有有限多个非零项的非负整数序列 $\{d_j\}$。——由定理 4，相交数
$\langle\tau_0^{d_0}\tau_1^{d_1}\cdots\tau_m^{d_m}\rangle_g$ 是体积多项式 $V_{g,n}(L)$
首项的系数。——

量子引力配分函数是
\[
\mathbf F(\lambda,\mathbf t)=\sum_{g=0}^{\infty}\lambda^{2g-2}F_g(\mathbf t).
\]

基于该函数用矩阵积分的一种实现，Witten 猜测配分函数应当满足 Korteweg--deVries（KdV）
方程的两个形式。Kontsevich 利用模空间 $\Mgn$ 的一个胞腔分解给出了该猜想的证明，[Kon92]。
胞腔由带状图（ribbon graph）／粗图（fat graph）枚举。Kontsevich 把相交数编码进三价
带状图的枚举。随后他用 Feynman 图方法与矩阵 Airy 积分建立了 Witten 的诸猜想。

相交数的两个基本关系是：对 $n>0$ 且 $\sum_i\alpha_i=3g-2+n>0$，有
\[
\text{string 方程}\qquad
\langle\tau_0\tau_{\alpha_1}\cdots\tau_{\alpha_n}\rangle_g
=\sum_{\alpha_i\neq 0}\langle\tau_{\alpha_1}\cdots\tau_{\alpha_{i-1}}\cdots\tau_{\alpha_n}\rangle_g,
\]

% [原文 p.32]
而对 $n\ge 0$ 且 $\sum_i\alpha_i=3g-3+n\ge 0$，有
\[
\text{dilaton 方程}\qquad
\langle\tau_1\tau_{\alpha_1}\cdots\tau_{\alpha_n}\rangle_g
=(2g-2+n)\langle\tau_{\alpha_1}\cdots\tau_{\alpha_n}\rangle_g.
\]

第一个方程对应新增一个穿刺点而不在乘积中附带 $\psi$ 因子，第二个方程对应新增一个
带单个 $\psi$ 因子的穿刺点。

与令 $V_{0,3}(L)=1$ 类似，带三个穿刺点的球面的相交符号被归一化为 $\langle\tau_0^3\rangle_0=1$。
一般的亏格 0 公式是
\[
\langle\tau_{\alpha_1}\cdots\tau_{\alpha_n}\rangle_0
=\binom{n-3}{\alpha_1!\cdots\alpha_n!},
\]
右端是 $n-3$ 的多重系数。亏格 0 的 string 方程不过是 Pascal 的多重系数邻项关系。

对亏格 1，定理 10 与 WP kappa 方程（定理 3）给出
$V(L)=\int_{\cM_{1,1}}2\pi^2\kappa_1+\frac{L^2}{2}\psi$。该公式与第 3 讲的计算
$V_{1,1}(L)=\frac{\pi^2}{12}+\frac{L^2}{48}$ 相结合（现在包含椭圆对合的 $\frac12$ 因子），
给出赋值
\[
\frac12\int_{\cM_{1,1}}\kappa_1=\langle\tau_1\rangle_1=\frac{1}{24}.
\]

% [Tutor Note: p.32 的 $V(L)=\int_{\cM_{1,1}}2\pi^2\kappa_1+\frac{L^2}{2}\psi$ 与紧随其后的 $\frac12\int_{\cM_{1,1}}\kappa_1=\langle\tau_1\rangle_1=\frac1{24}$ 要按椭圆对合的归一化一并阅读（$\Mgn$ 在 $(g,n)=(1,1)$ 时一般点的稳定化子是 $\Z_2$）。以 $V_{1,1}(L)=\frac{\pi^2}{12}+\frac{L^2}{48}$ 反推：常数项需 $2\pi^2\cdot\frac1{24}=\frac{\pi^2}{12}$，$L^2$ 系数需 $\frac12\cdot\frac1{24}=\frac1{48}$，两式恰好吻合，故前面的 $\frac12$ 与 $\int 2\pi^2\kappa_1$ 是配套的归一化；用标准值 $\langle\tau_1\rangle_1=\frac1{24}$ 核对时不要重复乘 $\frac12$。]

亏格 1 的一般赋值，由这一个赋值出发应用 string 与 dilaton 方程得到。Witten 猜想的一个
推论是：所有 $\langle\tau\rangle$ 相交数都可由初值 $\langle\tau_0^3\rangle_0=1$ 与
$\langle\tau_1\rangle_1=\frac{1}{24}$ 计算出来，所用的是下面描述的配分函数的 Virasoro 方程
$L_n(e^{\mathbf F})=0$。

在体积递推中，首项系数由首项系数得到——递推特殊化到首项系数，见
[Mir07b, Lemma 5.3]。我们现在简述把特殊化后的递推应用于配分函数的关系式，以及
Witten 猜想的一个解。

关系来自 Virasoro Lie 代数。Witt 子代数由微分算子 $\cL_n=-z^{n+1}\partial/\partial z$，
$n\ge-1$ 生成，其交换子为 $[\cL_n,\cL_m]=(n-m)\cL_{n+m}$。string 方程与 dilaton 方程
可以写成配分函数之指数 $e^{\mathbf F}$ 的线性齐次微分方程。string 方程的微分算子是
\[
L_{-1}=-\frac{\partial}{\partial t_0}+\frac{\lambda^{-2}}{2}t_0^2
+\sum_{j=0}^{\infty}t_{j+1}\frac{\partial}{\partial t_j},
\]
而 dilaton 方程的微分算子是
\[
L_0=-\frac{3}{2}\frac{\partial}{\partial t_1}
+\sum_{j=0}^{\infty}\frac{2j+1}{2}t_j\frac{\partial}{\partial t_j}+\frac{1}{16}.
\]
